\documentclass[11pt]{article}
\usepackage[margin=1in]{geometry}
\usepackage{amsmath,amssymb,amsthm,booktabs,microtype,hyperref}
\hypersetup{hidelinks,pdfauthor={T},pdftitle={Certifying Cumulative Harm Across Training Checkpoints}}
\newtheorem{proposition}{Proposition}
\title{Certifying Cumulative Harm Across Training Checkpoints}
\author{T}
\date{}
\begin{document}
\maketitle

\begin{abstract}
A later training stage can erase behavior learned during safety
tuning. We compute the nearest update satisfying one first-order
harm constraint, then show why that update alone can still raise
harm. A curvature term yields a sufficient finite-step condition.
For repeated updates along a fixed search line, we give the
sharp cumulative harm interval from a Lipschitz bound and
interval observations of the directional derivative. An exact
oracle checks the bounds and distinguishes preserved harm
allowance from a rule that resets its allowance after each step.
For checkpoints on a turning path, exact gradients and a global
smoothness bound give sharp endpoint values through shortest paths.
A source-aware search can certify a cap when a bounded number of
intermediate measurements may be faulty.
\end{abstract}

\section*{Summary and scope}
\paragraph{Question.} How much can a protected harm score increase across
a sequence of tentative updates?
\paragraph{Finding.} On a fixed search line, interval observations of a
Lipschitz directional derivative give exact endpoints for cumulative
change. Exact gradient vectors give a separate certificate that can
link nonadjacent checkpoints. One-sided projection or a renewed
allowance at every step can miss a breach.
\paragraph{Evidence and scope.} The interval is proved and
checked by a lattice oracle. A tuning team can certify each tentative
release against the same starting value. The guarantee applies to the
specified surrogate and a certified curvature bound.
\section{A protected behavioral objective}
Suppose $H(w)$ is a differentiable harm surrogate on a frozen
training slice. It might score whether a tool-use policy acts after
a veto, with the action validator separated from the reward model.
Let $a=\nabla H(w)\neq0$, and let $g$ be the update proposed by the
next training objective. Gradient surgery uses projections to
reduce gradient conflict in multitask learning \cite{pcgrad}.
Here a designated surrogate defines a one-sided guard for a
proposed update.
\begin{proposition}
The Euclidean-nearest update with $a^\top d\leq0$ is
\[
d=g-\frac{\max(0,a^\top g)}{\lVert a\rVert^2}a.
\]
If $\nabla H$ is $L$-Lipschitz on the update segment, then
$H(w+d)\leq H(w)+a^\top d+(L/2)\lVert d\rVert^2$.
\end{proposition}
\begin{proof}
When $g$ is feasible, its distance to itself is zero. Otherwise
write $g$ and any feasible vector in normal and tangential
coordinates. The tangential component can match $g$, while the
closest feasible normal coordinate is zero. The formula follows.
Integrating $\nabla H(w+td)^\top d$ for $0\leq t\leq1$ and
bounding its difference from $a^\top d$ by $Lt\lVert d\rVert^2$
gives the finite-step inequality.
\end{proof}

\section{Why a projected step can fail}
Take $H(x,y)=x+y^2$ at $(0,0)$, with $a=(1,0)$ and $g=(1,1)$.
Projection gives $(0,1)$, which satisfies the linear guard
with equality. Its harm rises from zero to one. If the permitted
slack is $m=H_{\max}-H(w)$, a sufficient update condition is
$a^\top d+(L/2)\lVert d\rVert^2\leq m$. With zero slack and no
first-order decrease, a finite step cannot be certified by this
inequality when $L>0$.

\section{A sharp cumulative bound on a fixed search line}
Let $w(x)=w_0+xv$ with $\lVert v\rVert=1$, and write
$g(x)=\frac{d}{dx}H(w(x))$. Suppose $g$ is $L$-Lipschitz
on $[0,T]$ and that certified observations restrict
$g(x_i)\in[\ell_i,u_i]$ at finitely many nodes, including
$0$ and $T$. A tentative checkpoint can supply the last
observation before an update is accepted.
Define the upper and lower extensions
\[
 U(x)=\min_i\{u_i+L|x-x_i|\},\qquad
 V(x)=\max_i\{\ell_i-L|x-x_i|\}.
\]
These are the interval-observation version of the extremal
McShane--Whitney envelopes for Lipschitz functions
\cite{petrakis}.

\begin{proposition}
The observations are feasible if and only if
$U(x_i)\ge\ell_i$ for every node. When feasible, the sharp
interval for $H(w(T))-H(w(0))$ is
\[
 \left[\int_0^T V(x)\,dx,\ \int_0^T U(x)\,dx\right].
\]
\end{proposition}
\begin{proof}
Every feasible $g$ lies between $V$ and $U$ by the Lipschitz
bound. Conversely a finite minimum or maximum of
$L$-Lipschitz functions is $L$-Lipschitz. At a node,
$U(x_i)\le u_i$, and the feasibility condition gives the
other inequality. That condition also states
$\ell_i\le u_j+L|x_i-x_j|$ for all node pairs. Reversing the
roles of the indices shows $V(x_i)\le u_i$, while the term
for node $i$ gives $V(x_i)\ge\ell_i$. Both envelopes are
admissible derivatives. Integrating each gives a witness
with the same initial value. Convex combinations of the two
derivatives attain intermediate changes.
\end{proof}

Between neighboring nodes at distance $d$, let $a$ and $b$
be the tightened values of $U$ at the left and right nodes.
All cones from the left have slope $L$ inside this interval
and all cones from the right have slope $-L$. Thus the
upper envelope at distance $t$ from the left is
$\min\{a+Lt,b+L(d-t)\}$. For $L>0$ its integral is
\[
 \frac{d(a+b)}{2}+\frac{Ld^2}{4}-\frac{(b-a)^2}{4L}.
\]
The formula follows by splitting the integral at
$t=(b-a+Ld)/(2L)$, which lies in $[0,d]$ by the Lipschitz
condition. For $L=0$ the feasible derivative is constant
between the nodes and the integral is $da$. Apply the same
upper calculation to $-g$ to get the lower endpoint.
With zero endpoint derivatives on a segment of length two
and $L=1$, the allowed increase can still reach one.

Keep the harm cap anchored to the protected reference
checkpoint. Certify the cumulative increase at every checkpoint
proposed for release. For example, with $H(w(x))=x$, two
increments of $3/5$ each fit under a freshly reset allowance
of one but their cumulative increase is $6/5$. If the search
turns to a different line, valid per-segment upper bounds can
be added from the same reference value. The resulting bound
need not be sharp on a path with changing directions.

\section{Linking checkpoints on a turning path}
Suppose the full gradients $g_i=\nabla H(x_i)$ are known at
finitely many checkpoints $x_i$. Assume $L>0$ and that the
gradient of $H$ is $L$-Lipschitz on the whole parameter space.
This stronger domain assumption lets the certificate use chords
between checkpoints as well as the proposed update segments.
Write $d_{ij}=x_j-x_i$ and set
\[
 c_{ij}=\frac12(g_i+g_j)^\top d_{ij}
 +\frac L4\lVert d_{ij}\rVert^2
 -\frac1{4L}\lVert g_j-g_i\rVert^2 .
\]

\begin{proposition}
The checkpoint gradients can come from a common globally
$L$-smooth surrogate if and only if the complete directed
graph with weights $c_{ij}$ has no negative cycle. In that
case the sharp interval for the change from checkpoint $0$
to checkpoint $T$ is
\[
 \bigl[-\operatorname{dist}(T,0),\ \operatorname{dist}(0,T)\bigr].
\]
\end{proposition}
\begin{proof}
Let $h_i$ be the unknown checkpoint values. Add the quadratic
$L\lVert x\rVert^2/2$ to $H$ to obtain a convex function
$F$ with $2L$-Lipschitz gradient. Conversely, subtracting that
quadratic from any convex $2L$-smooth function gives an
$L$-Lipschitz gradient. To see this, for a displacement $d$
and gradient difference $\Delta G$ of $F$, the smooth convex
gradient inequality gives
$\langle \Delta G,d\rangle\ge\lVert\Delta G\rVert^2/(2L)$.
Expanding $\lVert\Delta G-Ld\rVert^2$ bounds it by
$L^2\lVert d\rVert^2$.

The smooth convex interpolation theorem of Taylor, Hendrickx,
and Glineur \cite{taylor} requires
$F_i\ge F_j+\nabla F_j^\top(x_i-x_j)
 +\lVert\nabla F_i-\nabla F_j\rVert^2/(4L)$ for every
ordered pair. Substituting the shifted values and gradients
reduces this to $h_j-h_i\le c_{ij}$. Thus these difference
inequalities are necessary and sufficient for interpolation.

Summing around a cycle rules out a negative cycle. Conversely,
with no negative cycle, distances from any root $r$ give
feasible values $h_i=\operatorname{dist}(r,i)$ by the
triangle inequality for shortest paths. Summing along paths
from $0$ to $T$ bounds the increase above by
$\operatorname{dist}(0,T)$. The root $r=0$ attains this bound,
and $r=T$ attains the lower bound. Add a constant to match
the protected reference value. Interpolating between the
feasible value vectors gives all intermediate changes.
\end{proof}

For $L=1$, take $(0,0),(1,0),(0,1)$ in that order, with
zero gradient at each. The consecutive segment bounds sum to
$3/4$. The direct edge bounds the sharp increase by $1/4$.
Gradient consistency also requires more than the pairwise
Lipschitz check. At the same checkpoints, set the gradients
to $(0,0),(0,1),(-1,0)$. The two-edge cycles all have
weight zero, but the reverse triangle has weight $-1$.
These vectors cannot be gradients of the stipulated surrogate.

On 128 finite fixtures, graph closure agreed with direct
enumeration of simple paths and cycles. All 64 gradients
generated from smooth quadratics were feasible, and the
known values met every certified bound. The 64 unconstrained
gradient records were rejected. In 32 feasible fixtures the
graph improved on summing consecutive bounds. This is a
finite exact check of the graph calculation. It supplies
no measured global smoothness constant for a real model.

For gradient uncertainty sets, retain the vector $g_i$ as
a decision variable in each set and keep the original
difference inequalities. With convex uncertainty sets these
are convex constraints after moving the squared gradient
difference to the left side. This gives a finite formulation
for a sharp bound under the same smoothness model. The
shortest-path algorithm above assumes the gradients are fixed.

\subsection{A cap that tolerates faulty gradient sources}
Some intermediate gradient measurements can share a data pipeline.
Treat each such pipeline as one source and assume that at most
$k$ sources can be faulty. A faulty source contributes no trusted
gradients. The endpoint gradients and the global smoothness bound
remain trusted. We want the increase from $0$ to $T$ to be at most
$B$ for every admissible choice of faulty sources and every
surrogate that fits the remaining measurements.

\begin{proposition}
The following finite search either certifies the cap, finds a
compatible violation, or establishes that the measurements are
inconsistent within the fault budget. Begin with no sources removed.
At an inconsistent node, find a negative cycle and branch on its
removable sources. At a consistent node, let $P$ be a shortest
path from $0$ to $T$. If its weight exceeds $B$, report a violation.
Otherwise branch on the removable sources along $P$. Stop a branch
when the budget is exhausted or the witness contains no removable
source. Certify only if at least one consistent retained graph
was found and there is no violation.
\end{proposition}
\begin{proof}
Deleting a source removes its measurement vertices and leaves
all other weights unchanged. Any graph extending an inconsistent
node by further deletions must break the negative cycle before
it can become consistent. At a consistent node, the distance from
$0$ to $T$ is attained by a smooth surrogate, so a value above
$B$ is a valid counterexample. If the distance is at most $B$,
any graph that keeps the path $P$ still has an upper bound at
most $B$. A violating extension must therefore remove a source
on that path. These two witness rules preserve every possible
counterexample, and each branch consumes one source. If there
is no consistent node, all possible routes to consistency were
searched through the negative-cycle branches. This yields the
separate inconsistency decision.
\end{proof}

For example, take endpoints $(0,0),(2,0)$, $L=1$, and zero
measured gradients. Either intermediate checkpoint $(1,1/2)$
or $(1,-1/2)$ supplies a path with weight $5/8$. With one allowed
fault and distinct intermediate sources, at least one such path
survives. With both checkpoints owned by one source, losing that
source leaves the direct weight $1$, which is attainable. Two
measurements help only insofar as the fault model separates their
sources. A consistent full graph alone cannot establish tolerance
to deletions, since removing observations can increase the distance.

The witness search agreed with exhaustive source-subset enumeration
on 1,296 finite cap decisions. Of these, 597 were inconsistent,
353 certified and 346 admitted a violation. These counts describe
the verifier's grid. For a tuning workflow, record which checkpoint
measurements share an estimator, data slice or custody path before
setting the source budget. If the trusted assumptions fail, the
certificate supplies no basis for accepting the checkpoint.

\section{Integrating the guard into development}
During capability tuning, estimate protected gradients on a
declared replay slice and compare an ordinary update, a projected
update, and a projected update with backtracking. For every
tentative checkpoint, test the actual protected harm measures
before accepting it. Keep a separate final split from the repeated
line search. Log vetoed steps and stop when the objective cannot
improve within the allowed region. Multiple protected slices
require a common feasible set rather than a sequence of projections
assumed to preserve all earlier constraints.

\section{Local check and open test}
The closed form was checked against every feasible comparison on
the declared integer grid for 600 nonzero-normal and step pairs.
The proof establishes the continuous solution. The finite checks
recomputed 8,800 squared-distance inequalities. For the
fixed-line extension, a lattice dynamic program independently
recovered the exact endpoints for 352 compatible interval
histories and rejected all 160 incompatible ones. It agreed
on 2,112 harm-cap decisions. The sharp bound admitted 521
decisions that the left-endpoint-only certificate could not
admit on this fixed grid. A real model's
gradient noise and unknown curvature can invalidate a certificate,
and $H$ can miss harms. A next experiment compares checkpoint
retention under matched optimization budgets, including a case
where projection succeeds on replay and fails on a hidden slice.
\section{Sequential learning and the cost of retention}
We fitted a shared logistic scorer to a protected sign task,
then moved to a new task with a conflicting sign direction and a
free second feature. The protected arm retained at least 70\%
of the old sign margin and trained only the free direction in the
second stage. For the old feature support this implies a bound
on replay logistic loss. The matched unguarded arm updated both
directions. Across eight seeds and 3,200 held-out cases per task,
old-task accuracy was 0/3,200 unguarded and 3,200/3,200 protected.
New-task accuracy was 2,785/3,200 unguarded and 1,579/3,200
protected. All 288 proposed free-direction steps were accepted.
This example quantifies a severe retention cost in a deliberately
conflicting environment. It tests a restricted-direction guard,
with the earlier projection result governing a separate mechanism.

\paragraph{Review status.} Independent manuscript review and a language-model study remain open.

\begin{thebibliography}{9}
\bibitem{pcgrad} T. Yu et al. Gradient Surgery for Multi-Task Learning.
\url{https://arxiv.org/abs/2001.06782}.
\bibitem{petrakis} Iosif Petrakis. McShane--Whitney
Extensions in Constructive Analysis. \emph{Logical Methods
in Computer Science} 16(1), article 18, 2020.
\url{https://arxiv.org/abs/1804.06757}.
\bibitem{taylor} Adrien B. Taylor, Julien M. Hendrickx, and
Fran\c{c}ois Glineur. \emph{Smooth Strongly Convex Interpolation
and Exact Worst-Case Performance of First-Order Methods}.
Mathematical Programming 161, 307--345, 2017.
\url{https://arxiv.org/pdf/1502.05666}.
\end{thebibliography}
\end{document}
